\documentclass[12pt]{article}

\usepackage{fullpage}
\usepackage{multicol,multirow}
\usepackage{tabularx}
\usepackage{ulem}
\usepackage[utf8]{inputenc}
\usepackage{mathtools}
\usepackage{pgfplots}
\usepackage[russian]{babel}

\begin{document}

\section*{Лабораторная работа №\,8 по курсу дискрeтного анализа: Жадные алгоритмы}

Выполнил студент группы М8О-308Б-22 МАИ \textit{Немкова Анастасия}.

\subsection*{Условие}
На первой строке заданы два числа, N и p > 1, определяющие набор монет некоторой страны с номиналами $p^0$, $p^1$, ..., $p^{N-1}$. Нужно определить наименьшее количество монет, которое можно использовать для того, чтобы разменять заданную на второй строчке сумму денег M <=  $2^{32}$ - 1 и распечатать для каждого i-го номинала на i-ой строчке количество участвующих в размене монет.

\subsection*{Метод решения}

Жадный алгоритм решает задачу, принимая локально оптимальные решения на каждом этапе, с надеждой, что они приведут к глобально оптимальному решению.

В данной задаче, начиная с самой большой степени и двигаясь к меньшим, мы вычитаем из заданной суммы текущий номинал, пока это возможно, увеличивая счетчик для каждого номинала. Жадный алгоритм подходит для этого случая, так как номиналы образуют возрастающую последовательность степеней p. Для чисел, образующих геометрическую прогрессию, жадный выбор оптимален, так как каждый номинал можно представить как кратное более мелкого номинала. Поэтому, начиная с наибольшего номинала, мы всегда получаем минимальное количество вычитаний.

Если номиналы заданы произвольно и не имеют структуры геометрической прогрессии, то жадный подход может не привести к оптимальному решению. Например, для номиналов 1, 3, 4 значения m = 6, жадный алгоритм сначала выберет номинал 4 (1 раз), а затем дважды 1 (всего 3 слагаемых). Однако оптимальное решение — использовать два номинала 3, что даст 2 слагаемых.

Чтобы решить задачу для произвольного набора номиналов, можно использовать динамическое программирование. Суть алгоритма заключается  в том, что мы в массиве будем хранить минимальное количество номиналов, необходимых для получения суммы для каждого значения от 1 до m. В результате, массив будет содержать минимальное количество номиналов, необходимых для получения суммы m.

\subsection*{Описание программы}

Программа считывает три целых числа: n (количество номиналов), p (основание для степеней) и m (число, которое нужно разложить).

Создается вектор nominals, в который записываются степени числа p от $p^0$ до $p^{n - 1}$. Это делается с помощью цикла, который последовательно умножает предыдущее значение на p.

Далее инициализируется вектор counts, который будет хранить количество использованных номиналов для каждого значения. Затем, начиная с самого большого номинала, пытаемся вычесть его из m, пока m неменьше текущего номинала, увеличивая счетчик в counts для каждого использованного номинала.

Сложность: O(n⋅k), где n — количество степеней, а k — максимальное количество раз, которое может понадобиться для уменьшения m на номинал. В худшем случае, k может быть пропорционально m (если все номиналы маленькие).


\subsection*{Дневник отладки}

Были произведены попытки улучшить алгоритм, а также использовать динамическое программирование.

\subsection*{Тест производительности}

Для измерения производительсти сравним время работы данного алгоритма и наивного, в котором перебираем все возможные количества монет для каждого номинала и вычисляем сумму, создаваемую каждой комбинацией. Протестируем на входных данных, где p = 3, N = $k * 10$ и M = $k * 10^2$ (k = 1, ..., 5)
	
\begin{figure}[htbp]
    \centering
    \begin{tikzpicture}
        \begin{axis}[
            xlabel={Объём входных данных ($\times 10^2$)},
            ylabel={Время работы (ms)},
            grid=major,
            xmin=0.1, xmax=0.5,
            ymin=0, ymax=30,
            xtick={0.1, 0.2, 0.3, 0.4, 0.5},
            ytick={6, 12, 18, 24, 30},
            legend style={at={(0.5,-0.2)},anchor=north},
            legend columns=2,
            width=0.8\textwidth,
            height=0.5\textwidth,
            ]
            \addplot[color=blue,mark=*] coordinates {
                (0.1, 0.00071)
                (0.2, 0.00072)
                (0.3, 0.00075)
                (0.4, 0.0009)
                (0.5, 0.0013)

            };
            \addlegendentry{Мой алгоритм}
            \addplot[color=red,mark=square] coordinates {
                (0.1, 0.36)
                (0.2, 2.02)
                (0.3, 7.42)
                (0.4, 14.7)
                (0.5, 29.6)
            };
            \addlegendentry{Наивный}
        \end{axis}
    \end{tikzpicture}
    \label{fig:graph}
\end{figure}

Из графика видим, что жадный алгоритм имеет линейную сложность(O(n)), а наивный экспоненциальную (O($n * m^n$))

\newpage
\subsection*{Выводы}

В ходе данной лабораторной работы были разобраны жадные алгоритмы и их применение на примере размена денег с использованием монет различных номиналов. 

Основное отличие жадных алгоритмов от динамического программирования заключается в том, что жадные алгоритмы принимают решения на каждом шаге, основываясь на локально оптимальных выборах, в то время как динамическое программирование рассматривает более глобальный подход, учитывающий все возможные состояния для достижения оптимального решения. Это позволяет жадным алгоритмам часто работать быстрее и использовать меньше памяти, но они могут не находить оптимальное решение в задачах, где это требуется.

\end{document}




